% Part: normal-modal-logic
% Chapter: syntax-and-semantics
% Section: modal-validity

\documentclass[../../../include/open-logic-section]{subfiles}

\begin{document}

\olfileid{nml}{syn}{val}

\olsection{చెల్లుబాటుతనం}

\begin{explain}
  ప్రతి నమూనాలోనూ, అంటే ప్రతి నమూనాలోని ప్రతి లోకంలోనూ
  సత్యమయ్యే \tetoken{సూత్రాలు}{!!^{formula}s} ప్రత్యేకంగా ఆసక్తికరమైనవి.
  సాధ్యమైన లోకాల మధ్య ఉన్న ఏదో ఒక ప్రాప్యత సంబంధం నుంచి వచ్చే
  వ్యాఖ్యానాన్ని ఇక్కడ “నార్మల్” అంటాం. అలాంటి వ్యాఖ్యానం ఇచ్చినప్పుడు
  $\Box$, $\Diamond$లకు ఇచ్చే అర్థమేదైనా సత్యమయ్యే
  మోడల్ ప్రతిజ్ఞావాక్యాలను ఆ సూత్రాలు సూచిస్తాయి. అలాంటి
  \tetoken{సూత్రాలను}{!!{formula}s} \emph{చెల్లుబాటయ్యేవి}
  అంటాం. ఉదాహరణకు, $\Box(p \land q) \lif \Box p$
  చెల్లుబాటవుతుంది. అయితే, $\Box$కు అలెథిక్ (సత్య-అవశ్యకత)
  వ్యాఖ్యానం ఇచ్చినప్పుడు చెల్లుబాటవుతాయని అనిపించే కొన్ని
  \tetoken{సూత్రాలు}{!!{formula}s} ఉన్నాయి. ఉదాహరణకు,
  $\Box p \lif p$ చెల్లుబాటు కాదు. వేర్వేరు రకాల ప్రాప్యత
  సంబంధాల ద్వారా $\Box$, $\Diamond$లకు వేర్వేరు వ్యాఖ్యానాలు
  ఇవ్వగలగడం సంబంధాత్మక నమూనాల ప్రత్యేకతల్లో ఒకటి. అందువల్ల
  \emph{అన్ని} నమూనాల వర్గానికి సంబంధించీ,
  \emph{ఒక నిర్దిష్ట రకానికి చెందిన} అన్ని నమూనాల వర్గానికి
  సంబంధించీ చెల్లుబాటును నిర్వచించాలని ఇది సూచిస్తుంది. ఉదాహరణకు, ప్రతి లోకం తనకు తానే
  ప్రాప్యమయ్యే, అంటే $R$ స్వావర్తనమైన, అన్ని నమూనాల్లో
  $\Box p \lif p$ సత్యమని తేలుతుంది. నమూనాల వర్గాలకు
  సంబంధించి చెల్లుబాటును నిర్వచిస్తే ఈ విషయాన్ని సంక్షిప్తంగా
  చెప్పవచ్చు. స్వావర్తన నమూనాల వర్గంలో $\Box p \lif p$
  చెల్లుబాటవుతుంది.
\end{explain}

\begin{defn}
  నమూనాల వర్గం $\mClass{C}$లో \tetoken{సూత్రం}{!!^a{formula}}~$!A$
  \emph{చెల్లుబాటయ్యేది} $\mClass{C}$లోని ప్రతి నమూనాలోనూ అది
  సత్యమైనప్పుడు, అప్పుడే మాత్రమే (అంటే, $\mClass{C}$లోని ప్రతి నమూనాలోని
  ప్రతి లోకంలోనూ అది సత్యమైనప్పుడు, అప్పుడే మాత్రమే). $!A$ అనేది $\mClass{C}$లో
  చెల్లుబాటైతే $\mClass{C} \Entails !A$ అని రాస్తాం.
  \emph{అన్ని} నమూనాల వర్గంలో $!A$ చెల్లుబాటైతే
  $\Entails !A$ అని రాస్తాం.
\end{defn}

\begin{prop}\ollabel{prop:subset-class}
  $!A$ అనేది $\mClass{C}$లో చెల్లుబాటైతే,
  $\mClass{C}' \subseteq \mClass{C}$ అయ్యే ప్రతి
  ఉపవర్గంలోనూ చెల్లుబాటవుతుంది.
\end{prop}

\begin{prop}\ollabel{prop:Nec-rule}
  $!A$ చెల్లుబాటైతే $\Box!A$ కూడా చెల్లుబాటవుతుంది.
\end{prop}

\begin{proof}
  $\Entails !A$ అనుకుందాం. $\Entails \Box!A$ అని చూపడానికి,
  ఒక నమూనా $\mModel{M} = \tuple{W, R, V}$ను,
  అందులో ఏదో $w \in W$ను తీసుకుందాం. $Rww'$ అయితే, $!A$
  చెల్లుబాటవుతుంది కాబట్టి $\mSat{M}{!A}[w']$.
  అందువల్ల $\mSat{M}{\Box!A}[w]$ కూడా.
  $\mModel{M}$ను, $w$ను ఏవైనా తీసుకోవచ్చు కాబట్టి
  $\Entails \Box!A$.
\end{proof}

\begin{prob}
  కిందివి చెల్లుబాటవుతాయని చూపండి:
  \begin{enumerate}
  \item $\Box p \lif \Box (q \lif p)$;
  \item $\Box \lnot \lfalse$;
  \item $\Box p \lif (\Box q \lif \Box p)$.
  \end{enumerate}
\end{prob}

\begin{prob}
  $W = \{w\}$ అయ్యే $\mModel{M} = \tuple{W, R, V}$
  నమూనాల వర్గం $\mClass{C}$లో $!A \lif \Box!A$
  చెల్లుబాటవుతుందని చూపండి. అదేవిధంగా,
  $R = \emptyset$ అయ్యే $\mModel{M} = \tuple{W, R, V}$
  నమూనాల వర్గంలో $!B \lif \Box !A$, $\Diamond !A \lif !B$
  చెల్లుబాటవుతాయని చూపండి.
\end{prob}

\end{document}
